\chapter{Precondition Fuzzing}
\label{chap:precond}

The port's generator, like x86 \revizor{}'s, cannot produce the precondition violations behind
attacks such as PACMAN and TikTag. This chapter closes that gap with a methodology for
fuzzing instruction preconditions directly. The idea is to \emph{seal} the precondition sites
of a test case and fill each seal, per input, so that the precondition is deliberately
satisfied or violated on a speculative-only path, without ever faulting the victim
architecturally. \autoref{sec:precond-gap} states the gap, \autoref{sec:precond-sealing}
through \autoref{sec:precond-resolve} build the mechanism, \autoref{sec:precond-contract} fixes
the contract it tests against, \autoref{sec:precond-primitives} lists the primitives, and
\autoref{sec:precond-discussion} argues why the construction is sound.

\section{The generation gap}
\label{sec:precond-gap}

A precondition must hold along the architectural flow, but nothing constrains the speculative
paths, where a violation can leak without faulting the program. Random generation cannot target
this directly: whether an input makes a given instruction retire is an input-reachability
question, intractable in general. To stay safe for every input, the generator instruments every
potentially faulting instruction---masking memory addresses into the sandbox, avoiding
divide-by-zero, and so on---and so excludes the violation everywhere, transient paths included.
Precondition violations are excluded by construction. Reaching them requires a way to install a
violation that is guaranteed to be confined to a speculative-only path.

\section{Sealing a test case}
\label{sec:precond-sealing}

We walk a freshly generated test case statically. At each instruction whose operands we want to
fuzz we reserve a \emph{sealing}: a short, fixed-width area immediately preceding the instruction
of interest, initialized to \reg{UDF}, the permanently undefined A64 encoding. Sealing every site
yields the \emph{sealed test case}. Assembled once, it becomes an \emph{object test case}, a byte
template shared by all inputs. The \reg{UDF} filler guarantees that an unfilled sealing traps if
it is ever reached, so a sealing only ever executes once it has been deliberately filled.

Per input, we fill each sealing with instructions that enforce or violate its precondition just
before the guarded instruction---for example, materializing the correct \reg{MTE} tag or \reg{PAC}
signature to enforce the precondition, or an incorrect tag or signature to violate it. The fills
are carried by a \emph{relocation} section added to the input's \reg{.reif} file
(\autoref{sec:port-driving}). Just before the input runs, the loader applies the relocations and
patches the sealings dynamically. An input that satisfies every precondition is the \emph{genuine}
variant; an input with at least one unsatisfied, speculative-only sealing is a \emph{decoy}
variant. The two variants are then run and compared on hardware.

\section{The emulation table}
\label{sec:precond-emul}

Filling the genuine variant---and the architectural sealings of a decoy---requires values that
cannot be known statically, such as a pointer's \reg{PAC} signature or the allocation tag of the
granule an \reg{MTE} access touches. We give the contract executor (\autoref{sec:port-ce}) an
optional \emph{emulation table} whose rows bind each sealing to an \emph{op-emulator}. Each
op-emulator, in turn, dynamically patches the test case's instructions to satisfy the precondition
for the live register state, as if the variant had been genuine from the beginning. Because the CE
executes natively and models the extensions in software, it can compute these values even on a
host that does not implement the extension.

\section{Resolving a test case}
\label{sec:precond-resolve}

A sealing resolves differently for each input, because inputs take different paths and carry
different values. We recover what each sealing needs from a single CE trace of the object test
case, driven by the emulation table. The trace reports, per sealing, the minimal speculation depth
at which it is reached---depth zero means it retires, deeper means it is reached only
speculatively---and the value it must carry. This sidesteps the static reachability problem: for a
concrete input, the trace shows exactly which sites retire, so a violating fill can be placed only
where it is guaranteed to stay on a speculative-only path.

\section{The contract}
\label{sec:precond-contract}

The genuine and decoy variants run the \emph{exact same} input---the same data, \reg{PAC} keys,
and \reg{MTE} tags---and the same code outside the sealings. We test against CT-SEQ, which exposes
only the load and store addresses and the control flow of the in-order, non-speculative run. Every
violating sealing acts behind a speculative-only flow, so CT-SEQ requires it to leave the hardware
trace unchanged. Any divergence between the two variants is therefore a contract
violation~\cite{contracts,revizor}.

\section{Sealing primitives}
\label{sec:precond-primitives}

\autoref{tab:precond-sealings} lists the primitives. Each perturbs one precondition on a
speculative-only site, while the architectural path keeps the genuine fill. The Pointer
Authentication and Memory Tagging primitives are the automated analogues of PACMAN and TikTag: the
genuine fill installs a correct signature or tag, and the decoy an incorrect one, so a transiently
executed \reg{AUT*} or tagged access fails its check only on a squashed path. The canonicality and
alignment primitives, which \autoref{chap:canon} studies in depth, make the guarded address in
range and correctly aligned in the genuine variant and out of range or misaligned in the decoy.
The sign-extension primitive flips a high address bit so that the required sign extension no longer
holds; which bit is the effective sign bit depends on the Top Byte Ignore configuration (bit 55
when TBI is enabled, bit 63 otherwise), so the primitive flips the regime-correct bit.

\begin{table}[t]
\centering
\footnotesize
\setlength{\tabcolsep}{6pt}
\caption{Sealing primitives. Each perturbs one precondition on a speculative-only site, while the
architectural path keeps the genuine fill.}
\label{tab:precond-sealings}
\begin{tabular}{@{}l l l@{}}
\toprule
Precondition & genuine & decoy (speculative) \\
\midrule
PAC signature  & correct signature    & incorrect signature \\
MTE tag        & correct tag          & incorrect tag \\
Canonicality   & canonical address    & non-canonical address \\
Sign-extension & correct sign bits     & flipped sign bit \\
Branch target  & valid target         & non-canonical / misaligned \\
Page-table     & descriptor unchanged  & permissions flipped \\
\bottomrule
\end{tabular}
\end{table}

\paragraph{Page-table preconditions.}
Unlike the others, this primitive violates a system-configuration precondition rather than an
instruction operand. The sandbox maintains a pool of pages reachable only via speculative paths.
When a decoy variant executes, it randomly flips a controlled subset of a leaf page-table entry's
permission and attribute fields---the valid bit, the access flag, the access permissions, or the
memory-attribute index---while the genuine variant leaves every descriptor unchanged.

\section{Discussion}
\label{sec:precond-discussion}

Every retiring sealing satisfies its precondition by construction, so a violation only ever falls
on a speculative-only path and never crashes the test case. Each primitive also restores
architectural state immediately after the sealed access---it strips the forged \reg{PAC} bits,
retags to the real allocation tag, or undoes a pointer manipulation---so the genuine and decoy
variants re-converge immediately past the site, and their architectural difference is contained to
the single sealed instruction. The two variants are therefore not identical code: they differ
exactly in the perturbed sealings. A hardware-trace divergence between them is thus localized to
those sealings, and confirming a leak reduces to checking that the divergence reflects the
precondition being violated speculatively rather than an artifact of the perturbed instructions
themselves.
